\section{Analysis and Interpretation}
\label{sec:results}

The census establishes that an endpoint can remove a model before a score exists. This section asks
how far the surviving numbers move, and answers with the resolving power attached.

All numbers here come from the re-collection: an independent execution of the full
$4\times2\times2$ design on workspaces isolated per cell, comprising 5{,}809 valid measurements.
That this batch, rather than the original one, would carry the primary analysis was fixed in
writing before it produced a row (\S\ref{sec:threats}); the original collection is reported
alongside it as a replication, and the two agree on every criterion declared in advance.

% L'ORDINE DELLE SOTTOSEZIONI E' DELIBERATO. La decomposizione della varianza e l'inversione
% della calibrazione vengono prima della famiglia confermatoria: sono il risultato che regge, e
% mettere dieci test che non superano la correzione davanti a essi faceva arrivare il lettore alla
% parte debole per prima.
\subsection{Which lever binds is a property of the cell}
\label{sec:varianza}

The variance of a paired difference decomposes into within-binary noise, which shrinks with more runs
per binary, and heterogeneity between binaries, which shrinks only with more binaries. The two are not
interchangeable, and Table~\ref{tab:variance} shows the choice is not a property of the axis.

% Generata da analysis/tabella_principale.py --varianza su results/riraccolta.
\begin{table}[t]
\caption{Variance decomposition. $K_\infty$ is the number of binaries needed to resolve 3pp with unlimited runs per binary. $\dagger$: on these two contrasts the components are \textbf{poorly identified at the boundary} --- the estimated noise exceeds the observed variance, and the bootstrap clamps the residual to zero in 86\% of replications on T7 and 69\% on T8. No component is reported for them, in either direction: a clamped estimate is a boundary artefact, and printing a precise zero would report the constraint rather than the data. T1, T3 and T6 clamp in at most 1\% of replications, and are the rows the decomposition is read on.}
\label{tab:variance}
\centering
\small
\begin{tabular}{llrrr}
\toprule
& axis & SD & \% noise & $K_\infty$ \\
\midrule
T3 & transport & 0.2774 & 0.5\% & 668 \\
T6 & endpoint & 0.2189 & 9.1\% & 380 \\
T4 & transport & 0.1409 & 20.6\% & 137 \\
T1 & transport & 0.2002 & 46.1\% & 189 \\
T2 & transport & 0.0916 & 61.0\% & 28 \\
T5 & endpoint & 0.1292 & 93.1\% & 10 \\
T7 & endpoint & 0.0117 & \multicolumn{2}{c}{not decomposable$^\dagger$} \\
T8 & endpoint & 0.0585 & \multicolumn{2}{c}{not decomposable$^\dagger$} \\
\bottomrule
\end{tabular}
\end{table}


For T3 and T6 more runs buy nothing: 668 and 380 binaries against the 45 collected.

\textbf{These are estimates on 45 units, so we report their intervals.} Bootstrapping the binaries ---
the unit the contrast is paired on, 5{,}000 replications at a declared seed --- gives T3's noise share
as 0.5\% $[0.0, 1.4]$ and $K_\infty$ as 668 $[265, 1136]$, and T6's as 9.1\% $[4.1, 21.6]$ and 380
$[126, 744]$. The interval on $K_\infty$ is wide and its lower end is still six times the 45 binaries
collected, which is the operational point: the conclusion that more runs buy nothing does not depend on
the point estimate. Neither interval reaches the 50\% that would make the noise share
indistinguishable from half.

\textbf{Where the estimator is misspecified, we say so rather than offer two readings.} The rows marked
$\dagger$ have estimated noise exceeding the observed variance --- not a finding about those cells but a
variance model that does not fit them, since the decomposition assumes the two conditions' runs are
independent given the binary and clamps the residual at zero when that fails. The bootstrap makes the
failure countable: the residual is clamped in \textbf{86\%} of replications on T7, 69\% on T8 and 43\%
on T5, so for those contrasts the lower bound of $K_\infty$ is zero by construction rather than by
evidence, and we report no decomposition for T7 and T8. On T1, T3 and T6 --- the contrasts that carry
the finding --- clamping occurs in at most 1\% of replications.

\subsubsection*{What the metric's granularity does and does not explain}
\label{sec:granularita}

A decomposition of variance is only as trustworthy as the measurement it decomposes, and ours is coarse:
pass-rate over five unit tests gives a run \emph{six} values, 53\% of the 5{,}809 rows sit at an extreme and
33\% of the 720 binary-by-cell means land exactly on 0 or 1. A floor- and ceiling-heavy measurement is on its
own sufficient for a large between-binary spread, so it is a rival account of Table~\ref{tab:variance} and
not only of the inversion that follows.

\textbf{Two conditions bound the account, and both are stated so that a reader can check them.} It would
have to explain why the same metric on the same corpus yields 0.5\% noise here and 84\% in the collection
this design inherited its sizing from --- granularity is a property of the metric, and a property the two
share cannot by itself produce a difference between them. The second condition is narrower than the first and we state it as such: resampling binaries bounds the
sampling variance of the estimate, not a measurement bias, so the interval $[0.0, 1.4]$ on T3 does not by
itself exclude a floor-and-ceiling pattern uniform across the corpus --- such a pattern would reproduce in
every resample. The first condition carries the argument on its own, and neither is closed by these data
alone.

The operational consequence is in \S\ref{sec:proposizioni}: a design resting on a variance
decomposition should size the number of tests per task the way it sizes the number of tasks, and
report the share of runs at floor and ceiling beside the decomposition. We did the second and not
the first.

What T3 and T6 share is not the axis but the pairing of largest observed standard deviation with
smallest nominal $p$: where the effect looks largest it also varies most between binaries, which is
a weaker claim than a large average effect.

\subsection{RQ6 \emph{(exploratory, arose during execution)}: the power calibration
inherited from the previous chapter inverted}
\label{sec:rq6}

The pre-registration states, from the previous chapter's data, that for haiku ``84\% is
run-to-run noise and the floor is $K=13$''. On the same nominal contrast here, the noise share
is $0.5\%$ and the floor is $K=668$.

This is a reversed mechanism rather than a magnitude error, and it survives three independent changes:
of estimator, which widens it to 130.7\% noise and $K=0$; of apparatus and workspace isolation, where the
two collections give 0.1\%/705 and 0.5\%/668; and the other high-variance contrast, 9.3\%/487 against
9.1\%/380. The artifact carries each recomputation.

\textbf{The 84\% comes from the batch in which a validity flag failed, and the three robustness checks above
do not touch that.} Crashed rows were averaged as zeroes there (\S\ref{sec:method}), and the released
snapshot lets a reader reconstruct both computations. We claim no direction for that defect: independent
reconstructions disagree, and on that corpus the corrected estimator is misspecified in the sense of
\S\ref{sec:varianza}, so the quantity one would compare against is clamped rather than estimated. What the
reconstruction does establish is that the inherited 84\% is the \emph{defective} computation, which is why
the inversion is not read off the comparison: it is established inside this chapter by two collections of
our own, 0.1\%/705 and 0.5\%/668 on the contrast that carries it.

One caveat is removed by none of the three: the chapters differ in apparatus as well as estimator, and
exposure to the largest information leak of \S\ref{sec:threats} was asymmetric between arms in the earlier
chapter --- 100\% of native runs invoked the affected tool against 86.7\% of textual runs --- so the
comparison does not isolate the estimator. The inversion is larger than 13pp of differential exposure
makes plausible as a sole cause, but that is an order-of-magnitude argument, not an exclusion. What does
not depend on the comparison is this chapter's own decomposition: the 0.5\% noise share and the
668-binary floor are computed on our data alone and under one estimator, so the finding that the binding
lever is a property of the apparatus stands whether or not the earlier number is right.

\subsection{RQ1--RQ3 \emph{(confirmatory)}: a negative registered result, and the resolving
power that is the reportable quantity}

This subsection reports a \emph{negative} pre-registered result and reports it as such: hypotheses,
test family and analysis were fixed before the data existed, the family does not resolve an effect, and
the outcome is stated before the exploratory findings that did produce something. What it contributes
is not the absence of an effect but the \emph{resolution} at which the absence was established, which
differs by a factor of 23.7 between contrasts sharing one design.

Of the ten tests fixed in the pre-registration, \emph{none} passes its Holm threshold at
$m=10$~\citep{holm1979}. The smallest $p$ is $0.0191$ against a threshold of $0.0050$. No contrast places its
effect outside the $\pm 3$pp falsifier band with a confidence interval that excludes it.

\textbf{Part of that null follows from the family's fixed composition.} Two of the ten tests are not
independent of the other eight, so Holm at $m=10$ divides the smallest threshold by ten rather than
eight --- $0.0050$ instead of $0.0063$. T6's $p$ of $0.0191$ misses both, so no conclusion turns on the
choice here, though on a study whose smallest $p$ fell between them it would; we keep $m=10$ because the
pre-registration fixed it and because it is the conservative direction, not as the only defensible
count. The artifact carries both series.

That answers RQ1--RQ3 in the direction of ``not at this resolution'', and the resolution is the
reportable quantity: an evaluation that cannot separate 4pp from zero is not free to call two transports
equivalent. The outcome does not rest on the $t$'s normality assumption either ---
Appendix~\ref{sec:appendice-permutazione} repeats the family under a distribution-free sign-flip test and
no test passes Holm there.

We report the family outcome together with the achieved power, never one without the other, because
the three contrasts with nominal $p<0.05$ are also the three least powered --- and reporting power at
all is the exception here: 103 software engineering experiments were found systematically below
accepted norms~\citep{dyba2006power}, a decade after the concern was raised as ``missing and
misunderstood''~\citep{miller1997power}, and documented again for language-model
evaluation~\citep{card2020power}.

% Generata da analysis/tabella_principale.py su results/riraccolta — non modificare a mano.
\begin{table}[t]
\caption{The ten pre-registered tests, on the re-collection (isolated workspaces). $\delta$ is the paired difference per binary, over 45 binaries of 8 runs each. \textbf{Sign convention}: for transport contrasts $\delta = \text{textual} - \text{native}$, so a negative value means the textual protocol scores lower; for endpoint contrasts $\delta = \text{Bedrock} - \text{Databricks}$. MDE is the effect this contrast could have resolved at 80\% power given its own observed standard deviation, in percentage points --- the column exists because it varies by a factor of 23.7 across contrasts of one design. Power is computed against the pre-registered MDE of 4.87pp using the \emph{observed} standard deviation and the non-central $t$ --- the distribution the test actually uses. A normal approximation would report 0.0--1.9 points higher, which in a paper reporting its own under-powering would err in the flattering direction. $^\ddagger$The $p$ column is the \emph{exact} Student series, not the approximation the frozen script emits: the approximation understates $p$ in eight cases of eight --- for T6, 0.0155 against 0.0191 --- always in the direction that makes an effect look more significant. Both series are in the artifact, the distribution-free series is in Appendix~\ref{sec:appendice-permutazione}, and no test passes its Holm threshold under any of the three. $^{\P}$T5 is the one contrast whose $p$ rests on $K{=}44$ while its $\delta$ and interval rest on $K{=}45$: one binary returned seven valid runs of eight on one endpoint, the frozen script averages it and the exact-$p$ script requires a full complement. We print both rather than harmonise them, because harmonising would mean choosing a rule after seeing what it does to a pre-registered number.}
\label{tab:tests}
\centering
\scriptsize
\setlength{\tabcolsep}{2.2pt}
\begin{tabular}{llrrrrr}
\toprule
& contrast & $\delta$ & 95\% CI & $p$$^\ddagger$ & power & MDE \\
\midrule
T6 & llama,\, endp. & $-7.94$ & $[-14.5,-1.4]$ & 0.0191 & 30.9\% & 9.14 \\
T3 & haiku,\, transp. & $-9.83$ & $[-18.2,-1.5]$ & 0.0218 & 21.0\% & 11.59 \\
T1 & gpt-oss,\, transp. & $-6.56$ & $[-12.6,-0.5]$ & 0.0334 & 35.8\% & 8.36 \\
T9 & heterogeneity & --- & --- & 0.1011 & --- & --- \\
T2 & llama,\, transp. & $+2.22$ & $[-0.5,+5.0]$ & 0.1107 & 93.7\% & 3.82 \\
T5$^{\P}$ & gpt-oss,\, endp. & $+3.26$ & $[-0.6,+7.1]$ & 0.1159 & 69.6\% & 5.40 \\
T7 & haiku,\, endp. & $+0.22$ & $[-0.1,+0.6]$ & 0.2095 & 100.0\% & 0.49 \\
T8 & sonnet,\, endp. & $-0.94$ & $[-2.7,+0.8]$ & 0.2844 & 100.0\% & 2.44 \\
T10 & interaction & $-4.12$ & $[-10.6,+2.3]^{\S}$ & 0.3448 & --- & --- \\
T4 & sonnet,\, transp. & $-1.00$ & $[-5.2,+3.2]$ & 0.6363 & 62.1\% & 5.88 \\
\bottomrule
\end{tabular}
\end{table}


\begin{figure}[t]
\centering
\includegraphics[width=0.78\textwidth]{fig1-effetti-potenza.pdf}
\caption{The ten contrasts, ordered by nominal $p$. Filled and heavier lines are the three
with $p<0.05$ uncorrected; none survives Holm correction. They are also the three widest
intervals in the study --- the signature of a winner's curse, not of a larger effect. Dotted
lines mark the pre-registered $\pm3$pp falsifier band.}
\label{fig:effects}
\end{figure}

Table~\ref{tab:tests} and Fig.~\ref{fig:effects} are the whole confirmatory result. T3 reproduces the
transport difference that motivated this chapter at $-9.83$pp --- different apparatus, isolated
workspaces, three information leaks closed (\S\ref{sec:threats}) --- and T6 shows $7.94$pp between two
clouds serving the same nominal model at fixed transport. Neither survives correction for multiplicity,
and both sit below a third of the power the design was sized for. All three are true together.

$^{\S}$ The frozen script prints no interval for T10; this one is computed outside it with the \emph{model}
as the unit of replication ($n=4$, $t_3$). Treating the 180 binary-level pairs as independent would give a
narrower interval and a wrong one: they are four groups of 45.

\subsection{RQ1--RQ3 \emph{(confirmatory)}: the resolving power is per contrast, not a
single band}

The pre-registration sized the design with a pooled standard deviation of $0.1062$ and a
conservative $0.1167$, both calibrated on the previous chapter. The observed standard
deviations of the paired differences span $0.0117$ to $0.2774$ --- a ratio of 23.7 --- so a
single minimum detectable effect describes the design rather than the data. Computed per
contrast, the MDE ranges from $0.49$pp (haiku across endpoints) to $11.59$pp (haiku across
transports), and \emph{five of eight contrasts exceed} the 4.87pp the design assumed.

A null result here therefore reads: no contrast places its effect outside the band its own variance
permits. It does \emph{not} read as equivalence --- $K=45$ does not power an equivalence test at 3pp,
which would need 99--119 binaries, and this was declared before the data existed.

\subsection{RQ5 \emph{(exploratory, arose during execution)}: evidence against call
batching as the explanation}
\label{sec:rq5}

The pre-registered transport contrast moves two things at once on the models that batch: the format of
the call and the number a turn admits --- haiku issues 1.438 calls per turn against sonnet's 1.333. A
declared confound is not a bounded one, so we collected the arm that separates them: native transport
forced to \emph{one call per turn}, all else held fixed. It is exploratory, does not enter the family of
ten, and its stopping criterion was fixed before it ran.

\begin{table}[t]
\centering
\caption{Constraining the native transport to one call per turn. Paired on binaries, $K=45$,
eight runs per binary. The arm is exploratory and does not enter the corrected family.}
\label{tab:batching}
\footnotesize
\setlength{\tabcolsep}{4pt}
\begin{tabular}{@{}lrrrc@{}}
\toprule
model & full & constr. & diff. & 95\% CI \\
\midrule
haiku-4-5 & 0.747 & 0.754 & $+0.8$pp & $[-2.6, +4.2]$ \\
sonnet-4-5 & 0.793 & 0.809 & $+1.6$pp & $[-1.3, +4.4]$ \\
\bottomrule
\end{tabular}
\end{table}

Removing batching does not lower the score: both differences are \emph{positive} and both intervals
contain zero (Table~\ref{tab:batching}), against a transport effect of $-9.83$pp on the same model and
endpoint. The direction is what the arm supports --- a mechanism cannot explain a decrease it does not
produce --- and that inference needs only the sign of the two differences, which is why it survives the
arm's missed validity band (\S\ref{sec:rq5b}): the band governs how much the constraint bit, not which
way the score moved.

\emph{The magnitude does not follow, and we do not claim it.} ``Not within an order of magnitude''
would require the arm to have removed batching at the rate the band specified, and it did not. What
remains is narrower and still useful: the loss of call batching is not a plausible \emph{principal}
account of T3, and no reading of this arm establishes that the remainder is the protocol and nothing
else.

Two further cautions belong to the arm and are in the artifact: that the constraint does not remove
batching alone, so what it bounds is an upper limit rather than an estimate; and that the arm covers
one infrastructure and the two models that batch at all, since on the other two the rate is exactly
1.000 calls per turn and there is nothing to remove.

\subsection{RQ5, second finding \emph{(exploratory)}: one refusal reconfigures the model
for the rest of the run}
\label{sec:rq5b}

The arm produced a result we were not looking for, and it is sharp. Across 1{,}084 trajectories
\emph{every} first turn issues more than one call and is constrained; after the first refusal the rate
collapses to 0.2\%, eight turns out of 4{,}457. The model does not learn the constraint from the tool
schema, which does not encode it --- it learns it from one rejection, and the lesson holds for the
remainder of the trajectory.

\textbf{The arm falls outside its pre-declared validity band, and that governs what it licenses.} The
band required 33--38\% of turns to carry a constrained call, estimated from the unconstrained batching
rate; the observed figure is 19.7\%, outside it and not marginally. Under the rule fixed in
\S\ref{sec:design}, the ablation therefore does not license a statement about the \emph{size} of the
batching effect and we make none. Why the band was wrong is itself a measurement, and it is sharper than a missed target: the band assumed a
batching rate constant along the trajectory, and the collapse above fixes the achievable share at
essentially one constrained turn per trajectory. Of the 5{,}541 turns that carry a call, 1{,}084 are first
turns and every one of them is constrained; eight more are constrained later, for 1{,}092 in total, or
\textbf{19.7\%} --- which is the structural ceiling of 19.6\% plus the eight turns that still batched after
a refusal. The band's 33--38\% was never reachable under this constraint, so its premise was wrong rather
than the data.

The two statements are not interchangeable. \emph{The arm is not evidence about magnitude} --- that
reading requires the band it missed. \emph{It is evidence about direction}: both constrained cells score
no lower than unconstrained, which is incompatible with the loss of batching explaining a 9.83pp drop
whatever its size, since a mechanism cannot explain a decrease it does not produce. A sign survives a
failed magnitude check; an effect size does not. It also explains an asymmetry in the main comparison:
under the textual protocol the option to batch never exists, while under constrained native it exists
for one turn and the model abandons it. A transport is not only a format a call travels in --- it is a
channel that teaches.

\subsection{What differs between two clouds serving one model identity}
\label{sec:runtime}

T6 measures 7.94pp between two platforms serving the same model identifier, and \S\ref{sec:threats} states
that the weights cannot be verified identical --- which says nothing about \emph{how} the two runtimes
differ. Four quantities in the raw rows are determined by a serving stack rather than a model, and they do
not concord.

\begin{table}[t]
\centering
\caption{T6's two cells, native transport, 360 runs each. Median with first--third quartile; the
ratio is Bedrock over Databricks. Turn budget and corpus are identical by construction.}
\label{tab:runtime}
\footnotesize
\setlength{\tabcolsep}{4pt}
\begin{tabular}{@{}lrrr@{}}
\toprule
& Databricks & Bedrock & ratio \\
\midrule
input tokens & 20{,}562 & 8{,}853 & 0.43 \\
latency per run (s) & 13.4 & 8.6 & 0.64 \\
output tokens & 448 & 303 & 0.68 \\
candidate characters & 598 & 544 & 0.91 \\
turns & 12.0 & 12.0 & 1.00 \\
output tokens per second & 33.2 & 35.1 & 1.06 \\
\bottomrule
\end{tabular}
\end{table}

The two platforms hand the model \textbf{2.3 times} the input tokens for the same task, corpus and twelve
turns, while emitting output at within 6\% of the same rate. Throughput agreeing while input diverges
points away from hardware and towards what each stack \emph{sends} --- prompt scaffolding, tool-schema
serialisation, history truncation --- and which of the three is not visible from outside. This changes the
status of T6's confound rather than its size: declared before, documented now, with something in the two
stacks differing by a factor of 2.3 on the quantity the preceding chapter identified as a free parameter
in its own right~\citep{companion2026baseline}.

\subsection{Temperature zero is not determinism, and not equally so}

Two properties of the apparatus are usually assumed and turn out to be measurable. \emph{First,
temperature zero is not determinism.} All cells ran at temperature $0.0$, and across the eight runs of the
same binary in the same cell haiku returns an identical pass-rate on 91--100\% of binaries against
gpt-oss's 20--24\%, with 56\% of binaries over the study yielding eight identical runs
(Fig.~\ref{fig:determinism}).

\begin{figure}[t]
\centering
\includegraphics[width=0.78\textwidth]{fig2-determinismo.pdf}
\caption{Determinism at temperature $0.0$: the share of binaries whose eight runs return an
identical pass-rate, per model, with one point per cell. The spread within a model is the
distance between its clouds and transports --- the same nominal model is not equally stable
on different infrastructures.}
\label{fig:determinism}
\end{figure}

Second, \emph{the same nominal model is not equally stable on different clouds}, and on 45 binaries the
sampling supports that claim in exactly one of eight comparisons: llama on the native transport returns
identical scores on 24.4\% of binaries at one endpoint $[14.2, 38.7]$ and 66.7\% at the other $[52.1,
78.6]$, Wilson bands not touching. The other seven overlap and are not used as evidence. One clean case is
what these data support, on the same model and axis where T6 shows the largest endpoint effect in the
study, and the original collection isolates the same comparison.

\subsection{What the metric buys, measured}

Pass-rate has a floor reachable without reconstructing anything, and we measured it rather than
assuming it: the 53 runs that called \emph{no} tool average $0.1132$ against $0.6378$ for the 5{,}756
that called at least one. The floor exists and does not dominate --- the metric is a lower bound with a
measured floor, and the artifact carries the distribution.
\subsection{Declared covariates}

Two confounds were measured before the confirmatory numbers were read and neither is corrected for, since
correcting a pre-registered design after the fact is the freedom pre-registration exists to remove.
\emph{Calls per turn} is the one the ablation arm of \S\ref{sec:rq5} isolates, and it binds on two models
of four (\S\ref{sec:method}). \emph{Material acquired} is the second: replaying what each transport
actually requested, the native arm asks for more functions per binary in \emph{all eight} cells, most
strongly on llama at one endpoint ($-1.16$ functions, SD $0.90$), against six of eight originally. We do
not order that per-cell agreement against the size of the endpoint effect --- with four models the ordering
is indistinguishable from noise --- and the artifact carries both covariates in full.

\subsection{Scope}

The effect sizes are circumscribed twice over and both bounds are part of the result. They come from
\emph{one} family of tasks --- decompiled C binaries scored by unit tests --- on four models and two clouds,
so T3's $-9.83$pp is a measurement on this corpus rather than a quantity transferable elsewhere: a reader
who needs a number for their own setting should take the method and the variance decomposition, not the
point estimate. And every statement above concerns \emph{this} text protocol --- one call per turn, one line
of the response --- not textual tool calling in general, about which a leaderboard reporting both modes
finds the textual variant \emph{ahead} by 6pp on one model, against the direction that seeded this work.
